%==========================================================================
\section{Calibration inputs and fitting method}
\label{sec:method}

\paragraph{SG0 field scaling.} Revision~F divides the SG0 field, and hence $\BL_0$, by $f=7$. Only the
dose link~\eqref{eq:link} changes:
\begin{equation}
\D_{c,0}(I_0)=\D_{c,1}\,\frac{I_0}{f\,I_\star},\qquad fI_\star=7\times1.03869~\text{A}=7.2708~\text{A}.
\end{equation}
Because $\BL_0$ enters only through $B_0\ell_0^{\rm eff}$, this is equivalent to an effective SG0 length
$\ell_0^{\rm eff}=0.040/7=0.571$~cm at the measured field. $f$ is a \emph{fixed calibration input}; the
measured case $f=1$ is reported for comparison and Section~\ref{sec:fscan} maps the dependence on $f$.

\paragraph{Fixed inputs.} $\beta_0=0.031895$~T\,A$^{-1}$, $B_1=0.018931$~T, $\ell_0^{\rm eff}=0.040$~m,
$\ell_1^{\rm eff}=0.070$~m, $G_1=10.779$~T\,m$^{-1}$, $v_a=600.5$~m\,s$^{-1}$,
$|\gamma_e|=1.76086\times10^{11}$~rad\,s$^{-1}$\,T$^{-1}$, $K_{\rm pix}=6.47$~px, $D_0=13$~px. Unless
stated otherwise, $\DIR$ is solved from the baseline equation~\eqref{eq:dir}.

\paragraph{Objective.} Each data set is fitted by minimizing
\begin{equation}
\chi^2=\sum_{j=1}^{15}\left(\frac{\Delta D_j^{\rm model}-\Delta D_j^{\rm data}}{\sigma_j}\right)^2
\end{equation}
with the per-point uncertainties of the data file. Uncertainties on fitted parameters are profile
$\Delta\chi^2=1$ intervals (one-parameter fits) or curvature errors $\sqrt{2/\chi''}$. The model is
evaluated with Gauss--Legendre quadrature of \eqref{eq:zex}--\eqref{eq:Zpar}, and $\DIR$ by bisection
of~\eqref{eq:dir}.

\paragraph{Free parameters.} The parallel model has one free parameter, $\D_{c,1}$, equivalent to
$k_i=\D_{c,1}v_a/(|\gamma_e|\BL_1)$. Pure CQD adds none for the antiparallel set, so its antiparallel
curve is a prediction once $k_i$ is fixed. The phenomenological retained-label term adds one
parameter, $\kappa$ or $\RLC$.

%==========================================================================
\section{Recommended fit at \texorpdfstring{$f=7$}{f=7}: parallel set first, then \texorpdfstring{$k_i$}{ki} borrowed}
\label{sec:recommended}

\paragraph{Step 1: parallel set, pure CQD.} Minimizing $\chi^2_\pp(\D_{c,1})$ gives
\begin{equation}
\D_{c,1}=1.037,\qquad k_i=\frac{\D_{c,1}v_a}{|\gamma_e|\BL_1}=(2.67^{+0.60}_{-0.48})\times10^{-6},\qquad
\chi^2_\pp=2.83~(15~\text{points}),
\end{equation}
with $\alpha_0=\D_{c,1}/(fI_\star)=0.143$~A$^{-1}$, $\D_{c,0}(3.64~\text{A})=0.52$, $\Z_{0,\rm ex}=0.603$
and $\DIR=2.83$ ($L_d=0.198$~m, inferred). Fixing instead the reported $L_d=0.403$~m ($\DIR=5.757$)
gives $k_i=(2.59^{+0.51}_{-0.43})\times10^{-6}$ with $\chi^2_\pp=4.10$.

\paragraph{Step 2: antiparallel set with $k_i$ fixed.} Pure CQD now has no free parameter. Its strict
mirror-and-rebranch prediction rises to $+0.82$~px at 3.64~A, whereas the data fall to $-2.01$~px:
$\chi^2_\pa=109$. Fitting only the retained-label weight,
\begin{equation}
\Z^\pa=(1-w_{\rm ret})\Z^\pa_{\rm geom}+w_{\rm ret}\Z^\pa_{\rm ret},\qquad
w_{\rm ret}=\kappa\,\D_{c,0}\quad\text{or}\quad w_{\rm ret}=1-e^{-\RLC\,\D_{c,0}},
\end{equation}
gives $\kappa=0.249\pm0.024$ ($\chi^2_\pa=0.62$) or $\RLC=0.264\pm0.027$ ($\chi^2_\pa=0.80$). Both
forms vanish at $\D_{c,0}=0$ and so preserve the common baseline.

\begin{center}
\fbox{\parbox{0.92\textwidth}{\centering\textbf{Key result ($f=7$ fixed, $\DIR$ from the 13-px baseline, 2026-07-06 data)}\\[4pt]
Parallel, pure CQD: $\;k_i=(2.67^{+0.60}_{-0.48})\times10^{-6}$ ($\D_{c,1}=1.037$), $\;\chi^2=2.83$ for 15 points.\\[2pt]
Antiparallel, $k_i$ fixed: pure CQD $\chi^2=109$; with $w_{\rm ret}=\kappa\D_{c,0}$,
$\kappa=0.249\pm0.024$, $\chi^2=0.62$ for 15 points.\\[2pt]
Retained fraction: $w_{\rm ret}=\kappa\alpha_0I_0\approx0.036~\text{A}^{-1}\times I_0$, i.e.\ 13\% at 3.64~A.}}
\end{center}

\begin{figure}[h]
\centering
\begin{tikzpicture}
\begin{axis}[width=0.86\textwidth,height=0.52\textwidth,xlabel={SG0 current $I_0$ (A)},
  ylabel={$\Delta D$ (px)},xmin=-0.1,xmax=3.8,ymin=-3.0,ymax=6.6,grid=major,grid style={gray!20},
  legend pos=north west,legend style={font=\footnotesize,draw=none,fill=white,fill opacity=0.85,text opacity=1},legend cell align=left]
\addplot[parcol,thick] table[x=I0_A,y=dD_parallel_px,col sep=comma]{fit_curves_f7_20260706.csv};
\addlegendentry{parallel: pure CQD fit, $k_i=2.67\times10^{-6}$}
\addplot[antcol,thick] table[x=I0_A,y=dD_antiparallel_px,col sep=comma]{fit_curves_f7_20260706.csv};
\addlegendentry{antiparallel: same $k_i$, $w_{\rm ret}=\kappa\D_{c,0}$}
\addplot[gray,densely dotted,thick] table[x=I0_A,y=dD_antiparallel_pureCQD_px,col sep=comma]{fit_curves_f7_20260706.csv};
\addlegendentry{antiparallel: same $k_i$, pure CQD}
\addplot[only marks,mark=diamond*,mark size=3pt,parcol,error bars/.cd,y dir=both,y explicit]
  table[x=I0_A,y=dD_parallel_px,y error=sigma_parallel_px,col sep=comma]{spin_resolved_sep_20260706.csv};
\addlegendentry{parallel data}
\addplot[only marks,mark=square*,mark size=2pt,antcol,error bars/.cd,y dir=both,y explicit]
  table[x=I0_A,y=dD_antiparallel_px,y error=sigma_antiparallel_px,col sep=comma]{spin_resolved_sep_20260706.csv};
\addlegendentry{antiparallel data}
\addplot[gray,dashed,forget plot] coordinates{(-0.1,0)(3.8,0)};
\end{axis}
\end{tikzpicture}
\caption{Recommended reduction at $f=7$. The parallel curve is a one-parameter pure CQD fit; its $k_i$
is used unchanged for the antiparallel set, where only the retained-label weight is fitted. Curves are
tabulated in \texttt{fit\_curves\_f7\_20260706.csv}.}
\label{fig:recommended}
\end{figure}

\begin{table}[h]
\centering
\caption{Data and fit values at the measured currents ($\Delta D$ in px). ``$f=7$'' is the recommended
fit; ``$f=1$'' is the parallel fit with the measured SG0 field. The last column is the pure CQD
antiparallel prediction at the $f=7$ value of $k_i$.}
\label{tab:points}
\footnotesize\setlength{\tabcolsep}{5pt}
\begin{tabular}{lcccccc}
\toprule
 & \multicolumn{3}{c}{parallel} & \multicolumn{3}{c}{antiparallel}\\
\cmidrule(lr){2-4}\cmidrule(lr){5-7}
$I_0$ (A) & data & fit $f=7$ & fit $f=1$ & data & fit $f=7$ & pure CQD\\
\midrule
0.000 & $-0.047\pm0.50$ & 0.000 & 0.000 & $0.049\pm0.48$ & 0.000 & 0.000 \\
0.005 & $-0.045\pm0.50$ & 0.007 & 0.007 & $0.050\pm0.48$ & $-$0.002 & 0.000 \\
0.012 & $-0.043\pm0.50$ & 0.018 & 0.017 & $0.052\pm0.48$ & $-$0.005 & 0.000 \\
0.043 & $-0.030\pm0.51$ & 0.064 & 0.060 & $0.058\pm0.49$ & $-$0.019 & 0.000 \\
0.125 & $0.028\pm0.51$ & 0.185 & 0.173 & $0.063\pm0.50$ & $-$0.057 & 0.001 \\
0.431 & $0.442\pm0.50$ & 0.628 & 0.594 & $-0.016\pm0.52$ & $-$0.197 & 0.012 \\
0.937 & $1.343\pm0.47$ & 1.325 & 1.282 & $-0.293\pm0.51$ & $-$0.436 & 0.056 \\
1.404 & $2.034\pm0.47$ & 1.930 & 1.909 & $-0.551\pm0.50$ & $-$0.664 & 0.126 \\
1.403 & $2.033\pm0.47$ & 1.929 & 1.908 & $-0.551\pm0.50$ & $-$0.664 & 0.126 \\
1.867 & $2.539\pm0.50$ & 2.493 & 2.522 & $-0.780\pm0.51$ & $-$0.900 & 0.222 \\
2.327 & $3.014\pm0.53$ & 3.018 & 3.124 & $-1.033\pm0.52$ & $-$1.146 & 0.343 \\
2.780 & $3.608\pm0.52$ & 3.500 & 3.707 & $-1.357\pm0.53$ & $-$1.399 & 0.485 \\
3.221 & $4.296\pm0.49$ & 3.939 & 4.267 & $-1.729\pm0.52$ & $-$1.659 & 0.645 \\
3.642 & $4.775\pm0.47$ & 4.329 & 4.793 & $-2.011\pm0.51$ & $-$1.920 & 0.815 \\
3.598 & $4.748\pm0.47$ & 4.290 & 4.738 & $-1.991\pm0.51$ & $-$1.892 & 0.797 \\

\bottomrule
\end{tabular}
\end{table}

\paragraph{Residuals.} With $f=7$ the parallel curve bends over: the model lies 0.45~px below the
3.64~A point (about $1\sigma$) and 0.15--0.19~px above the data at 0.125 and 0.431~A. With $f=1$ the
parallel curve is nearly straight and all residuals are below 0.16~px ($\chi^2=0.48$). The
low-current points lie systematically below both models (the data start more slowly than linearly);
the effect is a third of $\sigma$ and is not significant with the present uncertainties.

\paragraph{Goodness of fit.} All $\chi^2$ values are well below the number of points (reduced
$\chi^2\approx0.04$--0.2). The quoted uncertainties are therefore conservative relative to the scatter,
or the points are correlated (for example through a common baseline subtraction, or if the closely
spaced low-current points are not independent measurements). The $\Delta\chi^2=1$ intervals should be
read with this in mind.

\subsection{Interpretation}
\begin{itemize}
\item \emph{$k_i$ is set by the slope of the parallel curve.} To leading order
$\Delta D^\pp/13\simeq2\D_{c,0}=2\alpha_0I_0$ (Section~\ref{sec:leading}), so the parallel data measure
$\alpha_0$, the SG0 dose per ampere. Converting $\alpha_0$ into $k_i$ uses the calibration, above all
$f$.
\item \emph{The factor 7 does not propagate as a factor 7.} With $f=7$ one needs $\D_{c,1}\approx1$,
where SG1 alone already aligns most of the branch ($\Z_0$ rises from 0.35 to 0.60). The relative effect
of a given SG0 pre-alignment is then smaller, and a larger SG0 dose is needed: $\alpha_0$ rises from
0.057 to 0.143~A$^{-1}$, and $k_i$ by a factor 17.5 rather than 7.
\item \emph{A small retained fraction explains the antiparallel decrease.} Retained-label atoms enter
SG1 counter-aligned and push strongly in the wrong direction, so 13\% at 3.64~A suffices. This is a
one-parameter phenomenological description, not a CQD prediction; whether it is physical must come
from a model of the field-reversal region (transit time $t_f$, adiabaticity, nuclear-angle motion).
\end{itemize}

\subsection{Results with the measured SG0 field (\texorpdfstring{$f=1$}{f=1})}
The same two-step procedure with $f=1$ gives $\D_{c,1}=0.0592$, $k_i=(1.52\pm0.09)\times10^{-7}$
($\chi^2_\pp=0.48$), $\DIR=5.20$ ($L_d=0.364$~m). For the antiparallel set the strict rule gives
$\chi^2=128$, and the retained-label weight gives $\kappa=0.452\pm0.040$ ($\chi^2=0.80$) or
$\RLC=0.470\pm0.044$ ($\chi^2=0.98$), i.e.\ $w_{\rm ret}(3.64~\text{A})=0.094$. Fixing $L_d=0.403$~m
leaves $k_i$ unchanged to four digits. These values agree with the digitized-data analysis
($1.53\times10^{-7}$, $\RLC=0.466$).

%==========================================================================
\section{Comparison of fitting strategies}
\label{sec:strategies}

\begin{table}[h]
\centering
\caption{Fitting strategies applied to the 2026-07-06 data. $\chi^2$ is for 15 points per set.}
\label{tab:strategies}
\footnotesize\setlength{\tabcolsep}{4pt}
\begin{tabular}{p{3.5cm}p{2.0cm}p{3.6cm}p{5.2cm}}
\toprule
Strategy & Free & $k_i$ ($f=1$ / $f=7$) & Verdict\\
\midrule
1. Both sets together (Figure-12 method) & $\D_{c,1}$, $\RLC$ & $1.52\times10^{-7}$ / $2.68\times10^{-6}$ ($\chi^2=1.46$ / 3.63, 30 pts) &
Works, but the non-CQD term $\RLC$ takes part in the fit that defines $k_i$. Gives the same $k_i$ as
strategy~2 within 0.5\%.\\[3pt]
2. Parallel only, then $k_i$ borrowed & $\D_{c,1}$; then $\kappa$ (or $\RLC$) & $1.52\times10^{-7}$ / $2.67\times10^{-6}$ ($\chi^2_\pp=0.48$ / 2.83) &
\textbf{Recommended.} $k_i$ comes from pure CQD on the data pure CQD can fit; the antiparallel set
then tests the reversal physics.\\[3pt]
3a. Antiparallel only, pure CQD & $\D_{c,1}$ & none &
Fails: the best it can do is $k_i\to0$, i.e.\ no change ($\chi^2\approx58$--60).\\[3pt]
3b. Antiparallel only, with retained label & $\D_{c,1}$, $\RLC$ & undetermined &
$\RLC$ compensates any $k_i$; $\chi^2$ stays near 1 (Figure~\ref{fig:strategies}, left).\\[3pt]
4. $f$ free as well & $\D_{c,1}$, ($\RLC$), $f$ & undetermined in $1.5\times10^{-7}$--$3\times10^{-6}$ &
Only $\alpha_0$ is measured; $\chi^2$ is flat for $f\lesssim6$ (Section~\ref{sec:fscan}).\\
\bottomrule
\end{tabular}
\end{table}

\begin{figure}[h]
\centering
\begin{tikzpicture}
\begin{axis}[width=0.62\textwidth,height=0.4\textwidth,xlabel={$k_i$},ylabel={$\chi^2_\pa$ (15 points)},
  xmode=log,ymode=log,xmin=5e-9,xmax=1e-5,ymin=0.1,ymax=4000,grid=major,grid style={gray!20},
  legend pos=north west,legend style={font=\scriptsize,draw=none,fill=white,fill opacity=0.85,text opacity=1},legend cell align=left]
\addplot[black,thick,densely dotted,mark=*,mark size=1.2pt] table[x=ki,y=chi2_f1_pure,col sep=comma]{plot_strategy3.csv};
\addlegendentry{$f=1$, pure CQD}
\addplot[black,thick,mark=*,mark size=1.2pt] table[x=ki,y=chi2_f1_rlc,col sep=comma]{plot_strategy3.csv};
\addlegendentry{$f=1$, with $\RLC$}
\addplot[parcol,thick,densely dotted,mark=square*,mark size=1.2pt] table[x=ki,y=chi2_f7_pure,col sep=comma]{plot_strategy3.csv};
\addlegendentry{$f=7$, pure CQD}
\addplot[parcol,thick,mark=square*,mark size=1.2pt] table[x=ki,y=chi2_f7_rlc,col sep=comma]{plot_strategy3.csv};
\addlegendentry{$f=7$, with $\RLC$}
\end{axis}
\end{tikzpicture}
\caption{Strategy 3: antiparallel set alone. Pure CQD never fits ($\chi^2\ge58$). With the retained-label
coefficient re-optimized at each $k_i$, $\chi^2$ stays near 1 over three decades for $f=7$, and for
$f=1$ it only bounds $k_i\lesssim1.5\times10^{-6}$.}
\label{fig:strategies}
\end{figure}

\paragraph{Why strategy 2 is preferred.} The parallel set is fitted by pure CQD with one parameter, so
$k_i$ is not influenced by any phenomenological term. Strategies 1 and 2 give the same $k_i$ within
0.5\%, which shows that the antiparallel data carry no independent information on $k_i$. Separating
the steps turns the antiparallel set into a clean test: with $k_i$ fixed, pure CQD has no freedom left
and fails ($\chi^2=109$), and the size of the retained-label fraction becomes a quantity to be
explained by a model of the field-reversal region.

%==========================================================================
\section{Dependence on the SG0 scaling \texorpdfstring{$f$}{f}}
\label{sec:fscan}

\begin{figure}[h]
\centering
\begin{tikzpicture}
\begin{groupplot}[group style={group size=2 by 1,horizontal sep=2cm},width=0.45\textwidth,height=0.37\textwidth,
  xlabel={SG0 field divisor $f$},grid=major,grid style={gray!20},xmin=0.5,xmax=10.5]
\nextgroupplot[ylabel={$k_i$},ymode=log,ymin=1e-7,ymax=1e-5]
\addplot[parcol,thick,mark=*,mark size=1.4pt] table[x=f,y=ki_par,col sep=comma]{plot_fscan.csv};
\nextgroupplot[ylabel={$\chi^2_\pp$ (15 points)},ymode=log,ymin=0.2,ymax=60]
\addplot[black,thick,mark=*,mark size=1.4pt] table[x=f,y=chi2_par,col sep=comma]{plot_fscan.csv};
\addplot[gray,dashed] coordinates{(0.5,23.7)(10.5,23.7)};
\end{groupplot}
\end{tikzpicture}
\caption{Parallel-only pure CQD fits as a function of the SG0 field divisor. Left: $k_i$. Right:
$\chi^2$; the dashed line is the 95\% limit for 14 degrees of freedom.}
\label{fig:fscan}
\end{figure}

The parallel fit is equally good for every $f\lesssim6$ ($\chi^2=0.46$--0.61), while $k_i$ grows from
$1.5\times10^{-7}$ to $1.8\times10^{-6}$. Beyond $f\approx6.5$ the fit degrades quickly:
$\chi^2=2.8$ at $f=7$, 11.9 at $f=8$ and 40 at $f=10$; the 95\% limit is crossed near $f\approx8.9$.
The data alone therefore measure $\alpha_0$, not $f$ and $k_i$ separately, and they cap $k_i$ at about
$3\times10^{-6}$ for any SG0 calibration. Between $f=6$ and $f=8$, $k_i$ changes from $1.8$ to
$2.9\times10^{-6}$, so the uncertainty of the factor 7 itself dominates the error on $k_i$.

\paragraph{Origin of the ceiling.} The initial slope of the parallel curve is
\begin{equation}
\frac{d\Delta D^\pp}{dI_0}\bigg|_0=\frac{13}{fI_\star}\;
\underbrace{\D_{c,1}\,\frac{\partial\Z^\pp/\partial\D_{c,0}}{\Z_0}}_{\equiv\,g(\D_{c,1})}.
\end{equation}
The capacity $g$ grows with $\D_{c,1}$ only until saturation sets in and peaks at $g\approx0.91$ near
$\D_{c,1}\approx2$ (Table~\ref{tab:g}). A measured slope therefore implies an upper bound on $f$.

\begin{table}[h]
\centering
\caption{Slope capacity $g(\D_{c,1})$ with $\DIR$ from the baseline.}
\label{tab:g}
\small\setlength{\tabcolsep}{4pt}
\begin{tabular}{lcccccccccc}
\toprule
$\D_{c,1}$ & 0.05 & 0.1 & 0.2 & 0.4 & 0.6 & 0.8 & 1.0 & 1.5 & 2.0 & 3.0\\
$\Z_0$ & 0.349 & 0.365 & 0.396 & 0.454 & 0.506 & 0.554 & 0.596 & 0.680 & 0.740 & 0.815\\
$g$ & 0.094 & 0.179 & 0.322 & 0.532 & 0.672 & 0.765 & 0.826 & 0.896 & 0.910 & 0.897\\
\bottomrule
\end{tabular}
\end{table}

\paragraph{Absolute geometry at $f=7$.} Solving $\DIR$ from the baseline gives $L_d=0.198$~m, which
contradicts both reported distances (0.320 and 0.403~m). Fixing $L_d=0.403$~m instead, the model
predicts a zero-current separation of 24.4~px rather than 13~px, which would require
$G_1\approx5.7$~T\,m$^{-1}$ instead of 10.8~T\,m$^{-1}$. The $f=7$ scaling brings $k_i$ close to the
numerical single-stage value ($2.51\times10^{-6}$) at the cost of moving the discrepancy into the
mechanical scale.

%==========================================================================
\section{Main caveats}
\begin{enumerate}
\item $f$, $\ell_0^{\rm eff}$ and $\ell_1^{\rm eff}$ are calibration inputs, not measured field integrals;
$k_i$ depends strongly on $f$ near $f=7$.
\item The absolute zero-current separation is taken as 13~px; it sets the vertical scale of every model
curve, since $\Delta D=13\,[\Z/\Z_0-1]$.
\item A single speed $v_a=600.5$~m\,s$^{-1}$ is used instead of the effusive speed distribution, and the
model computes branch centroids whereas the data are peak positions.
\item Reduced $\chi^2$ values of 0.04--0.2 indicate conservative or correlated uncertainties; a
generalized least-squares fit needs the covariance of the points.
\item The antiparallel fit depends on a phenomenological retained-label term that is not part of pure
CQD.
\end{enumerate}

%==========================================================================
\appendix
\section{Files}
All files are in the folder \texttt{fit\_20260706}.
\begin{description}
\item[\texttt{spin\_resolved\_sep\_20260706.csv}] \mbox{}\\ The data exported from the JLD2 file.
Columns: \texttt{I0\_A}, \texttt{dD\_parallel\_px},\\ \texttt{sigma\_parallel\_px},
\texttt{dD\_antiparallel\_px}, \texttt{sigma\_antiparallel\_px}.
\item[\texttt{fit\_curves\_f7\_20260706.csv}] \mbox{}\\ Recommended fit, $I_0=0$ to 3.70~A in steps of 0.01~A.
Columns: \texttt{I0\_A}, \texttt{dD\_parallel\_px}, \texttt{dD\_antiparallel\_px} (linear $w_{\rm ret}$),\\
\texttt{dD\_antiparallel\_exp\_px}, \texttt{dD\_antiparallel\_pureCQD\_px}, \texttt{Dc0}, \texttt{w\_ret}.
All values are changes relative to zero current (no 13-px offset).
\item[\texttt{fit\_curves\_f1\_20260706.csv}] \mbox{}\\ The same for $f=1$.
\item[\texttt{fit\_params\_f7\_20260706.csv}] \mbox{}\\ Parameters of the recommended fit.
\item[\texttt{results.json}] All fit results, including strategies 1--4 and the $f$ scan.
\item[\texttt{cqd\_model.js}, \texttt{run\_fits.js}] \mbox{}\\ Model and fits; \texttt{node run\_fits.js}
regenerates every number and CSV file (about 20~s). \texttt{make\_table.js} and
\texttt{build\_report.js} regenerate Table~\ref{tab:points} and this document.
\end{description}

\end{document}
